\section{Einleitung}
\label{sec:intro}

Die Bauindustrie befindet sich in einem tiefgreifenden digitalen Wandel, bei dem die Methode des Building Information Modeling (BIM) eine zentrale Rolle einnimmt. Das Ziel dieser Methode ist die Erfassung und Verwaltung aller relevanten Bauwerksdaten über den gesamten Lebenszyklus hinweg in einem digitalen Modell. Um eine Software-unabhängige Zusammenarbeit zu ermöglichen, setzen moderne Bauprojekte verstärkt auf den \textit{openBIM}-Ansatz, dessen technologisches Rückgrat der herstellerneutrale IFC-Standard (\textit{Industry Foundation Classes}) bildet.

Trotz der Standardisierungsbemühungen zeigt die Praxis, dass der Datenaustausch zwischen Planungsbeteiligten und ausführenden Bauunternehmen oft verlustbehaftet ist. Besonders in der zeitkritischen Angebotsphase stehen Unternehmen vor der Herausforderung, dass die erhaltenen Planungsmodelle zwar geometrisch detailliert, aber datentechnisch unvollständig für die erforderliche Mengenermittlung sind. Der daraus resultierende manuelle Nachbearbeitungsaufwand stellt ein erhebliches wirtschaftliches Risiko dar, da die Aufbereitung von Fremdmodellen in dieser frühen Projektphase oft ohne Erfolgsgarantie für den späteren Auftrag erfolgt.



An dieser Stelle setzt die vorliegende Arbeit an. Mit der Einführung der \textit{Information Delivery Specification} (IDS) durch buildingSMART steht ein modernes Format zur Verfügung, um Informationsanforderungen maschinenlesbar zu definieren. Bisher wird IDS primär als reines Prüfwerkzeug (Validierung) wahrgenommen. Das Ziel dieser Forschungsarbeit ist es jedoch, IDS als Grundlage für eine gezielte, teilautomatisierte Datenanreicherung (\textit{Enrichment}) von IFC-Modellen nutzbar zu machen.

Im Rahmen eines \textit{Design Science Research} Ansatzes wird ein technischer Prototyp entwickelt, der zeigt, wie Bauunternehmen externe Modelle effizient an ihre internen Datenstandards anpassen können[cite: 2, 3]. Dadurch soll nicht nur die Modellqualität messbar gesteigert, sondern auch der wirtschaftliche Aufwand in der Angebotskalkulation nachhaltig reduziert werden[cite: 2, 3].